% ============================================================
% DelegationChain paper — preamble
% Packages, settings, custom commands.
% ============================================================

% --- Document basics ---
\usepackage[a4paper, margin=2.5cm]{geometry}
\usepackage[T1]{fontenc}
\usepackage[utf8]{inputenc}
\usepackage{lmodern}
\usepackage{microtype}

% --- Math ---
\usepackage{amsmath, amsthm, amssymb}
\usepackage{mathtools}

% --- Tables and figures ---
\usepackage{booktabs}
\usepackage{array}
\usepackage{graphicx}
\usepackage{caption}
\captionsetup{font=small, labelfont=bf}

% --- Code listings ---
\usepackage{listings}
\usepackage{xcolor}
\definecolor{codebg}{HTML}{F7F7F7}
\definecolor{codeborder}{HTML}{D0D0D0}
\lstset{
    basicstyle=\ttfamily\small,
    backgroundcolor=\color{codebg},
    frame=single,
    rulecolor=\color{codeborder},
    breaklines=true,
    keepspaces=true,
    columns=fullflexible,
    xleftmargin=4pt,
    xrightmargin=4pt,
    aboveskip=8pt,
    belowskip=8pt,
}

% --- Cross-references ---
\usepackage{hyperref}
\hypersetup{
    colorlinks=true,
    linkcolor=black,
    urlcolor=blue!50!black,
    citecolor=blue!50!black,
    pdftitle={DelegationChain},
    pdfauthor={Daniyar Amangeldin},
}
\usepackage[capitalize]{cleveref}

% --- Bibliography ---
\usepackage[
    backend=biber,
    style=authoryear,
    sorting=nyt,
    natbib=true,
    maxcitenames=2,
    mincitenames=1,
    maxbibnames=99,
]{biblatex}
\addbibresource{references.bib}

% --- Theorem environments ---
\newtheorem{theorem}{Theorem}
\newtheorem{lemma}[theorem]{Lemma}
\newtheorem{proposition}{Proposition}
\theoremstyle{definition}
\newtheorem{definition}{Definition}
\theoremstyle{remark}
\newtheorem*{remark}{Remark}

% --- Spacing ---
\setlength{\parskip}{0.4em}
\setlength{\parindent}{1.2em}

% --- Custom commands ---
\newcommand{\G}[1]{\mathbb{G}_{#1}}
\newcommand{\Gt}{\mathbb{G}_T}
\newcommand{\Zq}{\mathbb{Z}_q}
\newcommand{\pair}[2]{e(#1, #2)}
\newcommand{\Sign}{\mathsf{Sign}}
\newcommand{\Verify}{\mathsf{Verify}}
\newcommand{\HashG}{\mathsf{HashToG}_2}
\newcommand{\Canon}{\mathsf{Canon}}
\newcommand{\sk}{\mathit{sk}}
\newcommand{\pk}{\mathit{pk}}

% Threat label macro
\newcommand{\threat}[2]{\textbf{T#1:} \emph{#2}}